\section{Detailed Related Work}
\label{app:related}

The closest literatures all study model response, but they attach different objects and semantics to that response. We organize the comparison around the object being explained and the question it answers.

\subsection{Gradient, path, and signed attribution}

Gradient explanations measure local sensitivity, while path methods accumulate sensitivity from a baseline to the input~\citep{simonyan2014deep,selvaraju2017gradcam,sundararajan2017axiomatic,smilkov2017smoothgrad,xu2020attribution,kapishnikov2021guided,you2026joint}. Layer-wise relevance propagation, DeepLIFT, FullGrad, and SHAP instead decompose a prediction into additive feature contributions under conservation, reference, or game-theoretic principles~\citep{bach2015lrp,shrikumar2017deeplift,srinivas2019fullgrad,lundberg2017unified}. Some of these methods preserve positive and negative contributions~\citep{dhurandhar2018pertinent}.

DECAF studies a different object. It does not distribute one prediction across input coordinates. It decomposes the scalar response of an already specified paired intervention. The clean endpoint supplies the orientation, and the endpoint-null branch separates fragility from positive or negative evidence. Spatial saliency and DECAF are therefore complementary: saliency indicates where a response is localized, whereas DECAF indicates what semantic role the paired response plays.

\subsection{Perturbation, counterfactual, and causal explanation}

Occlusion, Meaningful Perturbations, Extremal Perturbations, RISE, LIME, and removal-based Shapley methods construct interventions and summarize the resulting model changes~\citep{zeiler2014visualizing,fong2017meaningful,fong2019extremal,petsiuk2018rise,ribeiro2016why,covert2021removing}. Counterfactual and contrastive methods retrieve or synthesize changes that alter a decision or support a contrast~\citep{goyal2019counterfactual,chang2019counterfactual,dhurandhar2018pertinent,pmlr-v258-you25a,gu2026discover}. Causal attribution and concept-intervention methods interpret relevance through explicit interventions on inputs, representations, or human-interpretable concepts~\citep{chattopadhyay2019causal,janzing2020causal,goyal2020cace,kim2018tcav,koh2020conceptbottleneck}.

DECAF takes the paired intervention as input rather than searching for it. Counterfactual construction defines the factor and identification assumptions; the paired trajectory measures the response; endpoint orientation then decomposes that response into evidence, contradiction, and fragility. The same decomposition can therefore sit on top of different removal operators or counterfactual generators.

\subsection{Baselines, off-manifold effects, and faithfulness evaluation}

Baseline and path choices can substantially change attribution~\citep{sturmfels2020baselines,kindermans2019unreliability}, and perturbations can move samples away from the data manifold~\citep{frye2020manifold}. Sanity checks, retraining-based evaluations, infidelity measures, and debiased removal metrics test whether explanations reflect the model rather than artifacts of the evaluation protocol~\citep{adebayo2018sanity,hooker2019benchmark,yeh2019infidelity,rong2022consistent}. Other work documents label leakage, explanation fragility, and adversarial manipulation~\citep{jethani2023labelleakage,ghorbani2019interpretation,dombrowski2019explanations,slack2020fooling}.

DECAF does not claim path invariance. It keeps the intervention protocol explicit and asks which component changes when the protocol changes. This distinction matters in ImageNet-9: nested-patch reveal raises total response by roughly $80\%$, but the increase comes primarily from contradiction and fragility rather than evidence. DECAF addresses a semantic ambiguity that remains even when a paired response has been measured faithfully.

\subsection{Efficiency, black-box access, and benchmark context}

Black-box perturbation methods trade model access for repeated queries; RISE and dense occlusion can require many forward evaluations, while amortized methods such as FastSHAP add a separate explainer-training problem~\citep{petsiuk2018rise,jethani2022fastshap}. DECAF instead reuses the signed scores already computed by a paired perturbation curve. The routing step adds no model queries, requires no gradients or internal activations, and can be batched across examples, stages, factors, and models.

ImageNet-9 is a standard setting for studying background reliance and distribution shift~\citep{xiao2021noise,geirhos2020shortcut}. DECAF asks a finer question than whether a model responds to background: is that response endpoint-aligned evidence, endpoint-opposed contradiction, or sensitivity that appears only when the endpoint effect is null?
